\documentclass[12pt,a4paper]{report}
\usepackage[utf8]{inputenc}
\usepackage[french]{babel}
\usepackage[T1]{fontenc}
\usepackage{graphicx}
\usepackage{hyperref}
\usepackage{listings}
\usepackage{xcolor}
\usepackage{geometry}
\usepackage{titlesec}
\usepackage{fancyhdr}

\geometry{margin=2.5cm}

% Configuration des couleurs pour le code
\definecolor{codegreen}{rgb}{0,0.6,0}
\definecolor{codegray}{rgb}{0.5,0.5,0.5}
\definecolor{codepurple}{rgb}{0.58,0,0.82}
\definecolor{backcolour}{rgb}{0.95,0.95,0.92}

\lstdefinestyle{mystyle}{
    backgroundcolor=\color{backcolour},   
    commentstyle=\color{codegreen},
    keywordstyle=\color{magenta},
    numberstyle=\tiny\color{codegray},
    stringstyle=\color{codepurple},
    basicstyle=\ttfamily\footnotesize,
    breakatwhitespace=false,         
    breaklines=true,                 
    captionpos=b,                    
    keepspaces=true,                 
    numbers=left,                    
    numbersep=5pt,                  
    showspaces=false,                
    showstringspaces=false,
    showtabs=false,                  
    tabsize=2
}

\lstset{style=mystyle}

% Configuration des en-têtes et pieds de page
\pagestyle{fancy}
\fancyhf{}
\fancyhead[L]{SECURITY UNIVERSITY}
\fancyhead[R]{Rapport de Projet}
\fancyfoot[C]{\thepage}

\begin{document}

\begin{titlepage}
    \centering
    \vspace*{1cm}
    {\huge\bfseries SECURITY UNIVERSITY\par}
    \vspace{1cm}
    {\Large Rapport de Projet\par}
    \vspace{1.5cm}
    {\large\itshape Équipe de Développement\par}
    \vfill
    {\large \today\par}
\end{titlepage}

\tableofcontents
\newpage

\chapter{Introduction}
\section{Contexte du Projet}
SECURITY UNIVERSITY est une application de gestion universitaire sécurisée, conçue pour faciliter l'administration et la gestion des données académiques tout en maintenant des standards élevés de sécurité.

\section{Objectifs}
\begin{itemize}
    \item Développer une plateforme sécurisée de gestion universitaire
    \item Assurer la protection des données sensibles
    \item Fournir une interface utilisateur intuitive
    \item Permettre une gestion efficace des ressources académiques
\end{itemize}

\chapter{Architecture du Système}
\section{Technologies Utilisées}
\begin{itemize}
    \item Backend : PHP 7.4+
    \item Base de données : MySQL 5.7+
    \item Frontend : HTML5, CSS3, JavaScript
    \item Versioning : Git
\end{itemize}

\section{Structure du Projet}
\begin{lstlisting}[language=bash]
SECURITY_UNIVERSITY/
├── assets/
│   ├── css/
│   ├── js/
│   └── images/
├── config/
├── database/
├── functions/
├── includes/
├── modules/
└── templates/
\end{lstlisting}

\chapter{Fonctionnalités}
\section{Gestion des Utilisateurs}
\begin{itemize}
    \item Authentification sécurisée
    \item Gestion des rôles et permissions
    \item Profils utilisateurs
\end{itemize}

\section{Gestion Académique}
\begin{itemize}
    \item Gestion des étudiants
    \item Gestion des cours
    \item Système de notation
    \item Emplois du temps
\end{itemize}

\chapter{Sécurité}
\section{Mesures de Sécurité}
\begin{itemize}
    \item Authentification à deux facteurs
    \item Chiffrement des données sensibles
    \item Protection contre les injections SQL
    \item Gestion sécurisée des sessions
\end{itemize}

\section{Protection des Données}
\begin{itemize}
    \item Chiffrement des mots de passe
    \item Protection des données personnelles
    \item Journalisation des accès
\end{itemize}

\chapter{Installation et Configuration}
\section{Prérequis}
\begin{itemize}
    \item PHP 7.4 ou supérieur
    \item MySQL 5.7 ou supérieur
    \item Serveur Web (Apache/Nginx)
\end{itemize}

\section{Procédure d'Installation}
\begin{enumerate}
    \item Installation des dépendances
    \item Configuration de la base de données
    \item Configuration du serveur web
    \item Mise en place des permissions
\end{enumerate}

\chapter{Interface Utilisateur}
\section{Design et Ergonomie}
\begin{itemize}
    \item Interface responsive
    \item Navigation intuitive
    \item Accessibilité
\end{itemize}

\section{Composants Principaux}
\begin{itemize}
    \item Tableau de bord
    \item Formulaires de gestion
    \item Rapports et statistiques
\end{itemize}

\chapter{Conclusion}
\section{Bilan du Projet}
Le projet SECURITY UNIVERSITY répond aux besoins de gestion universitaire tout en assurant un haut niveau de sécurité des données.

\section{Perspectives}
\begin{itemize}
    \item Intégration de nouvelles fonctionnalités
    \item Amélioration continue de la sécurité
    \item Optimisation des performances
\end{itemize}

\appendix
\chapter{Annexes}
\section{Documentation Technique}
\subsection{Configuration du Serveur}
\begin{lstlisting}[language=apache]
# Configuration Apache
<VirtualHost *:80>
    ServerName security.university
    DocumentRoot /var/www/html/SECURITY_UNIVERSITY
    <Directory /var/www/html/SECURITY_UNIVERSITY>
        Options Indexes FollowSymLinks
        AllowOverride All
        Require all granted
    </Directory>
</VirtualHost>
\end{lstlisting}

\subsection{Configuration de la Base de Données}
\begin{lstlisting}[language=sql]
-- Structure de la base de données
CREATE TABLE users (
    id INT PRIMARY KEY AUTO_INCREMENT,
    username VARCHAR(50) UNIQUE,
    password VARCHAR(255),
    role VARCHAR(20),
    created_at TIMESTAMP DEFAULT CURRENT_TIMESTAMP
);
\end{lstlisting}

\end{document} 